\begin{tabular}{@{}p{0.16\linewidth}p{0.80\linewidth}@{}}
\toprule
项目 & 设置 \\
\midrule
海流数据 & HYCOM 全球海洋再分析（GLBy0.08，实验 93.0），水平分辨率 1/12°，取 10 米深度；23°--31°N、121°--129°E（台湾以东至东海，含黑潮）；2019-01-01 至 2024-08-29 每 3 天 12:00 时一张，共 689 张 \\
任务 & 每张快照抽 2 个：起终点相距 100--140 千米、离岸不少于 20 千米，两点中间设半径 15 千米的禁航区，可从左侧或右侧绕过；只取粗图上任务区平均流速不低于 0.3 米/秒的强流任务；开发 710 个（2019--2021 年），测试 623 个（2022--2024 年） \\
粗图与真实海流 & 航行器带 3 天前的 1° 粗分辨率海流图（下称粗图，格内平均）；真实海流为当天 1/12° 再分析海流场 \\
航行器 & 对水速度 1.0 米/秒；选定一侧后，按真实海流在该侧以最快方式航行 \\
航线求解 & 网格最短路：粗网格求解 8 千米网格、8 个方向；中网格求解 4 千米、16 个方向；细网格求解 2 千米、32 个方向（计算量之比 1 : 8 : 64）；顶流过强的边不可通行 \\
合理流场集合 & 粗图加上由流函数生成的无散误差模态的组合，各点误差大小与相关长度（76 千米）由开发年份标定，组合系数服从标准正态；半径 $r=1.5$（开发年份中包住真实左右航时差的任务占 93\%） \\
合理反例检验 & 在半径 $r$ 的球内对左右航时差做投影梯度上升（梯度沿最快航线累加）；两个方向、半径 1 与 2 各搜一次，其他半径线性插值，每个任务约 71 次中网格求解 \\
观测方式 & 原地停留测流（出发前 2 小时测 12 次）；沿途测流 30 千米或沿途测流 60 千米后可换边；绕行定点测流（去前方 40 千米内反例与粗图差别最大处测量）；上浮取当天 1° 图或上浮取当天 0.5° 图（每次 1 小时）；测流噪声 0.05 米/秒 \\
阈值判据 & 最大后悔值超过 8 小时（开发年份选定）时判定粗图不可信，触发超阈值时保守改选或超阈值时上浮取图，否则按粗图选边 \\
评价指标 & 额外航时 = 实际航时 $-$ 真实更快一侧的航时（小时）；配对差值的 95\% 置信区间按天整组自助重抽样 2000 次 \\
\bottomrule
\end{tabular}
